% Original analytic teaching diagram, 2026-11-09.
\definecolor{VB}{HTML}{7188A3}\definecolor{VBF}{HTML}{E5EAF0}
\definecolor{VG}{HTML}{829C8B}\definecolor{VGF}{HTML}{E8EEE9}
\definecolor{VY}{HTML}{B4A16F}\definecolor{VYF}{HTML}{F1ECDD}
\definecolor{VR}{HTML}{A7787D}\definecolor{VRF}{HTML}{F0E5E6}
\definecolor{VP}{HTML}{9588A5}\definecolor{VPF}{HTML}{ECE8F0}
\definecolor{VI}{HTML}{3B4252}\definecolor{VM}{HTML}{6B7280}\definecolor{VL}{HTML}{4C4D4F}
\begin{tikzpicture}[x=1mm,y=1mm,>=Stealth,font=\figurefont\fontsize{8}{10}\selectfont,
lab/.style={text=VM,font=\figurefont\fontsize{7.5}{9}\selectfont}]
% (a) Jacobian split into B and N
\draw[VG,fill=VGF,line width=.8pt] (0,6) rectangle (20,26);
\draw[VY,fill=VYF,line width=.8pt] (20,6) rectangle (56,26);
\node[text=VI] at (10,16) {$\bm{B}$};
\node[text=VI] at (38,16) {$\bm{N}$};
\draw[VL,line width=.7pt] (0,27.5)--(20,27.5);
\draw[VL,line width=.7pt] (20,27.5)--(56,27.5);
\foreach \a in {0,20,56}{\draw[VL,line width=.7pt] (\a,26.8)--(\a,28.2);}
\node[text=VI,anchor=south] at (10,28) {$\bm{y}\in\mathbb{R}^{p}$};
\node[text=VI,anchor=south] at (38,28) {$\bm{z}\in\mathbb{R}^{n-p}$};
\node[lab,anchor=north] at (10,5.2) {基本变量};
\node[lab,anchor=north] at (38,5.2) {非基本变量};
\node[text=VI,rotate=90,anchor=south] at (-1,16) {$p$ 个方程};
\node[text=VI] at (28,-6) {(a)};
% (b) crossing lines y^2=z^2: no basis at the origin
\begin{scope}[shift={(82,16)}]
\draw[VL,line width=.6pt,->] (-20,0)--(20,0) node[anchor=north east,lab] {$z$};
\draw[VL,line width=.6pt,->] (0,-14)--(0,15) node[anchor=north west,lab] {$y$};
\draw[VL,line width=1.6pt] (-14,-14)--(14,14);
\draw[VL,line width=1.6pt] (-14,14)--(14,-14);
\node[circle,fill=VR,inner sep=0pt,minimum size=2mm] at (0,0) {};
\node[text=VI] at (0,-22) {(b)};
\end{scope}
\end{tikzpicture}
